\begin{proof}[Proof of \cref{obj:lem:radius-sensitive-phiw}]
\begin{enumerate}
\item
Write \(\mathbb P_M^{\mathrm{obs}}\) for the observed trajectory law induced by \(M\), and let
\[
q_C:=\frac{C-1}{C}
\]
be the overlap-distance radius from \cref{obj:def:overlap-distance-radius}.  Define the unclipped partial-history average
\[
\widetilde\theta_k
:=
\frac{1}{T-k}\sum_{t=k}^{T-1} Z_t(k),
\qquad
\widehat\theta_k
=
\operatorname{clip}_{[-1,1]}(\widetilde\theta_k),
\]
with \(Z_t(k)\) as in \cref{obj:def:phiw-estimator}.  Since \(1\le T\) and \(k\le \lfloor T/2\rfloor\), one has \(k<T\), so \(T-k\ge 1\).  The conditions defining \(M\in\mathcal M_T(t_0,\zeta,C)\) in \cref{obj:def:model-class} give \(L=e^\zeta>1\) and \(0<\alpha=e^{-1/t_0}<1\).  The finite observation space makes each score \(Z_t(k)\) measurable.  Moreover \(L=e^\zeta\ge 1\), and the sequential-ignorability, policy-overlap, and bounded-reward restrictions in \cref{obj:def:model-class} imply the window envelope
\[
|Z_t(k)|\le L^{k+1}
\qquad \mathbb P_M^{\mathrm{obs}}\text{-a.s.},
\]
for every \(t\).  Thus each score is square-integrable under \(\mathbb P_M^{\mathrm{obs}}\), and so is \(\widetilde\theta_k\).  Moreover, \cref{obj:lem:target-value-unit-interval} applies to the present horizon and model class membership, giving
\[
\theta(K,e)\in[-1,1].
\] % lean: CausalSmith.Stat.PomdpLatentOverlapMinimax.radius_sensitive_phiw

\item
For every \(u\in\mathbb R\) and every \(\vartheta\in[-1,1]\),
\[
\bigl(\operatorname{clip}_{[-1,1]}(u)-\vartheta\bigr)^2
\le
(u-\vartheta)^2.
\]
Applying this pointwise with \(u=\widetilde\theta_k\) and \(\vartheta=\theta(K,e)\), and then using the usual bias--variance identity for the square-integrable random variable \(\widetilde\theta_k\), yields
\[
\mathbb E_M\!\left[
\{\widehat\theta_k-\theta(K,e)\}^2
\right]
\le
\operatorname{Var}_M(\widetilde\theta_k)
+
\left(
\mathbb E_M[\widetilde\theta_k]-\theta(K,e)
\right)^2 .
\] % lean: CausalSmith.Stat.PomdpLatentOverlapMinimax.sqRisk_clipUnit_le_variance_add_bias_sq

\item
We next bound the bias term.  Let \(d_b\) and \(d_e\) be the stationary laws of the behavior- and target-policy transition kernels, and let \(g_e\) be the target-policy reward regression appearing in \cref{obj:def:target-value}.  Define
\[
g_e^\circ(s)
:=
\begin{cases}
g_e(s), & d_b(s)>0,\\
0, & d_b(s)=0,
\end{cases}
\qquad
\Psi_k(s):=(P_e^k g_e^\circ)(s).
\]
On states with \(d_b(s)>0\), the horizon condition and the model-class restrictions give \(|g_e(s)|\le 1\); on states with \(d_b(s)=0\), the definition gives \(g_e^\circ(s)=0\).  Hence \(|g_e^\circ(s)|\le 1\) for every \(s\), and therefore
\[
\operatorname{osc}(g_e^\circ)
:=
\sup_{s,s'} |g_e^\circ(s)-g_e^\circ(s')|
\le 2.
\]
The kernel and policy restrictions in \cref{obj:def:model-class} make \(P_e\) a Markov transition kernel, and \cref{obj:ass:uniform-contraction} supplies
\[
\|\lambda P_e-\lambda' P_e\|_{\mathrm{TV}}
\le
\alpha\|\lambda-\lambda'\|_{\mathrm{TV}}
\]
for all probability laws \(\lambda,\lambda'\) on \(\mathcal S\).  Consequently, for any bounded function \(r:\mathcal S\to\mathbb R\) and any states \(s,s'\), the dual total-variation inequality gives
\[
\bigl|(P_e r)(s)-(P_e r)(s')\bigr|
\le
\operatorname{osc}(r)\,\|\delta_sP_e-\delta_{s'}P_e\|_{\mathrm{TV}}
\le
\alpha\operatorname{osc}(r).
\]
Iterating this one-step oscillation contraction yields
\[
\operatorname{osc}(P_e^k r)
\le
\alpha^k\operatorname{osc}(r).
\]
Applying the last display to \(r=g_e^\circ\) gives
\[
\operatorname{osc}(\Psi_k)
=
\operatorname{osc}(P_e^k g_e^\circ)
\le
2\alpha^k .
\]
By \cref{obj:lem:phiw-score-stationary-mean}, every score in the average defining \(\widetilde\theta_k\) has mean
\[
\mathbb E_M[Z_t(k)]
=
\sum_{s\in\mathcal S} d_b(s)(P_e^k g_e)(s),
\qquad t=k,\ldots,T-1.
\]
Therefore
\[
\mathbb E_M[\widetilde\theta_k]
=
\sum_{s\in\mathcal S} d_b(s)(P_e^k g_e)(s).
\]
For each state \(s\) with \(d_b(s)>0\), the functions \(g_e\) and \(g_e^\circ\) agree on the behavior support, so \cref{obj:lem:behavior-support-operator-congruence}, applied with the target policy and initial state \(s\), gives
\[
(P_e^k g_e)(s)=\Psi_k(s).
\]
For states with \(d_b(s)=0\), the common prefactor \(d_b(s)\) makes both weighted summands vanish, regardless of the two function values.  Hence
\[
\sum_s d_b(s)(P_e^k g_e)(s)
=
\sum_s d_b(s)\Psi_k(s).
\]
By \cref{obj:ass:latent-stationary-overlap}, included among the restrictions in \cref{obj:def:model-class}, \(d_e(s)\le C d_b(s)\).  Thus \(d_e(s)>0\) implies \(d_b(s)>0\), so
\[
\theta(K,e)
=
\sum_s d_e(s)g_e(s)
=
\sum_s d_e(s)g_e^\circ(s).
\]
By the meaning of the target-policy stationary law used in \cref{obj:def:target-value}, \(d_e\) is invariant for \(P_e\).  Thus, for every function \(r_{\mathrm{test}}:\mathcal S\to\mathbb R\) and every integer \(m\ge0\),
\[
\sum_s d_e(s)r_{\mathrm{test}}(s)
=
\sum_s d_e(s)(P_e^m r_{\mathrm{test}})(s).
\]
Applying this identity with \(r_{\mathrm{test}}=g_e^\circ\) and \(m=k\) yields
\[
\sum_s d_e(s)g_e^\circ(s)
=
\sum_s d_e(s)\Psi_k(s).
\]
Combining the preceding displays,
\[
\mathbb E_M[\widetilde\theta_k]-\theta(K,e)
=
\sum_s \{d_b(s)-d_e(s)\}\Psi_k(s).
\]
For any function with oscillation at most \(2\alpha^k\),
\[
\left|
\sum_s \{d_b(s)-d_e(s)\}\Psi_k(s)
\right|
\le
2\alpha^k\,\|d_b-d_e\|_{\mathrm{TV}}.
\]
It remains to prove the stationary-law distance subclaim
\[
\|d_b-d_e\|_{\mathrm{TV}}\le q_C.
\]
This subclaim uses the pointwise comparison \(d_e(s)\le C d_b(s)\) from \cref{obj:ass:latent-stationary-overlap}.  Throughout this step,
\[
\|\mu-\eta\|_{\mathrm{TV}}
:=
\frac12\sum_{s\in\mathcal S}|\mu(s)-\eta(s)|
=
\sup_{B\subseteq\mathcal S}|\mu(B)-\eta(B)|
\]
for probability laws on the finite joint-state space.  With this convention, any two probability laws have total-variation distance at most \(1\), since
\[
\frac12\sum_s |\mu(s)-\eta(s)|
\le
\frac12\sum_s\{\mu(s)+\eta(s)\}=1.
\]
If \(C=1\), then \(d_e(s)\le d_b(s)\) for every \(s\), and the two probability laws have equal total mass, so \(d_b=d_e\).  Hence
\[
\|d_b-d_e\|_{\mathrm{TV}}=0=q_C.
\]
If \(C>1\), define
\[
\eta_C(s):=\frac{C d_b(s)-d_e(s)}{C-1},\qquad s\in\mathcal S .
\]
The latent stationary overlap condition gives \(C d_b(s)-d_e(s)\ge0\) for every \(s\), and
\[
\sum_s \eta_C(s)
=
\frac{C\sum_s d_b(s)-\sum_s d_e(s)}{C-1}
=
1,
\]
so \(\eta_C\) is a probability law.  The identity
\[
d_b=\frac1C d_e+q_C\eta_C,
\qquad q_C=\frac{C-1}{C},
\]
then gives
\[
d_b-d_e
=
q_C(\eta_C-d_e).
\]
Therefore
\[
\|d_b-d_e\|_{\mathrm{TV}}
=
q_C\|\eta_C-d_e\|_{\mathrm{TV}}
\le q_C.
\]
Thus in all cases,
\[
\left|
\mathbb E_M[\widetilde\theta_k]-\theta(K,e)
\right|
\le
2q_C\alpha^k .
\] % lean: CausalSmith.Stat.PomdpLatentOverlapMinimax.phiw_bias_le

\item
The variance estimate supplied by \cref{obj:lem:phiw-raw-variance}, with \(L=e^\zeta\) and \(\alpha=e^{-1/t_0}\), gives
\[
\operatorname{Var}_M(\widetilde\theta_k)
\le
\frac{2}{T}\left\{
L^{k+1}\left(1+\frac{4}{L-1}\right)
+
\frac{4}{1-\alpha}
\right\}.
\] % lean: CausalSmith.Stat.PomdpLatentOverlapMinimax.variance_phiwRaw_le

\item
Since \(q_C\ge 0\) and \(\alpha\ge 0\),
\[
\left(
\mathbb E_M[\widetilde\theta_k]-\theta(K,e)
\right)^2
\le
(2q_C\alpha^k)^2
=
4q_C^2\alpha^{2k}.
\]
Substituting this and the variance bound into the clipped bias--variance inequality proves
\[
\mathbb E_M\!\left[\{\widehat\theta_k-\theta(K,e)\}^2\right]
\le
4q_C^2\alpha^{2k}
+\frac{2}{T}\left\{
L^{k+1}\left(1+\frac{4}{L-1}\right)
+\frac{4}{1-\alpha}
\right\}.
\] % lean: CausalSmith.Stat.PomdpLatentOverlapMinimax.radius_sensitive_phiw
\end{enumerate}
\end{proof}
